%% ---------------------------------------------------------------------------
%% intro.tex
%% Introducción
%% ---------------------------------------------------------------------------

\chapter{Introducción}
\label{chp:intro}

La seguridad en los sistemas embebidos ha adquirido un papel
fundamental debido al crecimiento acelerado del Internet de las Cosas (IoT), los
dispositivos móviles y las infraestructuras industriales interconectadas. En este nuevo
escenario tecnológico, millones de dispositivos operan de forma distribuida y, en muchos
casos, en entornos físicamente accesibles, lo que incrementa significativamente la
superficie de ataque \cite{shamsoshoara_2020_puf}. Como consecuencia, garantizar la autenticidad, integridad y
confiabilidad de cada dispositivo se convierte en un requisito esencial para prevenir
clonación, falsificación y accesos no autorizados.

Tradicionalmente, la seguridad criptográfica se fundamenta en el uso de claves secretas
almacenadas en memorias no volátiles. Sin embargo, este enfoque introduce una debilidad
estructural: la información sensible permanece físicamente almacenada dentro del
dispositivo, quedando expuesta a ataques invasivos y semi-invasivos capaces de extraer
las claves directamente del hardware \cite{hou_2019_lightweight}. Esta problemática es
especialmente crítica en plataformas embebidas de bajo costo y recursos limitados, donde
la incorporación de mecanismos avanzados de protección física resulta compleja o
económicamente inviable.

Ante esta limitación, las Funciones Físicamente No Clonables (PUFs, por sus siglas en inglés
\textit{Physical Unclonable Functions}) han emergido como una alternativa prometedora
para el establecimiento de raíces de confianza en hardware. Las PUFs explotan las
variaciones físicas inherentes al proceso de fabricación de los
circuitos integrados para generar respuestas únicas e irrepetibles, eliminando la
necesidad de almacenar claves criptográficas de forma permanente. Gracias a esta
propiedad, las PUFs permiten implementar mecanismos de autenticación, identificación de
dispositivos y generación dinámica de claves con un nivel elevado de resistencia frente
a ataques físicos \cite{shamsoshoara_2020_puf, handschuh_2011_hardware}.

Dentro de las múltiples arquitecturas existentes, la variante
PLPUF (\textit{Pseudo Linear Feedback Shift Register PUF}) representa una solución
particularmente atractiva para plataformas reconfigurables. A diferencia de otras
implementaciones basadas en retardos o elementos de memoria, el PLPUF emplea lógica
combinacional para generar respuestas multibit a partir de un único desafío, logrando
una implementación compacta, eficiente en área y altamente parametrizable
\cite{hori_2011_pseudolfsr}. Estas características lo convierten en un candidato idóneo
para su integración en sistemas embebidos modernos basados en FPGA.

En este contexto, el presente trabajo propone el diseño, implementación e integración
de un módulo PLPUF programable como periférico dentro de un
\textit{System-on-Chip} (SoC) basado en FPGA. El módulo se conecta mediante la interfaz
estándar AXI y es controlado por el procesador suave MicroBlaze~V sobre la plataforma
Basys3 con FPGA Artix-7. Este enfoque elimina la dependencia de almacenamiento persistente de claves,
fortaleciendo la seguridad a nivel físico, y provee una plataforma experimental
para evaluar esquemas de autenticación y generación dinámica de claves
en sistemas embebidos con recursos limitados.

De esta manera, el proyecto se posiciona en la intersección entre el diseño de sistemas
embebidos y la seguridad de hardware, contribuyendo al desarrollo de soluciones abiertas,
reconfigurables y reproducibles orientadas a fortalecer la seguridad desde el propio
nivel físico del dispositivo.

% ---------------------------------------------------------------------------
\section{Antecedentes del proyecto}

\subsection{Descripción de la institución}

El Instituto Tecnológico de Costa Rica (TEC) es una institución pública de educación
superior fundada en 1971 con el propósito de impulsar el desarrollo científico y
tecnológico del país. A lo largo de su trayectoria, el TEC se ha consolidado como un
referente nacional en formación profesional e investigación aplicada, destacándose por
su enfoque en innovación tecnológica, excelencia académica y compromiso con el desarrollo
sostenible \cite{TEC_general}.

El presente proyecto se desarrolla en la Escuela de Ingeniería en Computadores, unidad
académica orientada a la integración de conocimientos en Electrónica y Ciencias de la
Computación. Esta escuela promueve activamente la participación estudiantil en proyectos
de investigación y desarrollo en áreas como diseño digital, sistemas embebidos,
verificación de hardware, inteligencia artificial y seguridad informática.

\subsection{Área de conocimiento del proyecto}

El trabajo se enmarca dentro del área de Sistemas Embebidos, al involucrar el diseño e
implementación de un módulo hardware sobre una arquitectura reconfigurable (FPGA) y su
integración dentro de un System-on-Chip (SoC) con procesador y periféricos de
comunicación.

Paralelamente, el proyecto se vincula con la Seguridad de Hardware mediante la
implementación de una Función Física No Clonable (PUF) como raíz de confianza,
orientada a la generación de identificadores únicos, protección de datos y autenticación
de dispositivos sin requerir almacenamiento permanente de claves criptográficas.

\subsection{Trabajos similares}

La arquitectura Pseudo-LFSR PUF (PLPUF) fue propuesta originalmente por
Hori~et~al.~\cite{hori_2011_pseudolfsr}, quienes demostraron que una cadena
de inversores con retroalimentación tipo Fibonacci, equivalente a un LFSR
implementado de forma puramente combinacional, es capaz de generar respuestas
multibit compactas y eficientes en área sobre FPGAs Xilinx Spartan. Su
diseño alcanzó niveles competitivos de uniformidad y unicidad con un consumo
reducido de recursos lógicos. Más recientemente,
Alsharkawy~et~al.~\cite{alsharkawy_2024_plpufs} analizaron los riesgos
ocultos de la duración de activación en PLPUFs, evidenciando que una
selección inadecuada de este parámetro puede comprometer la calidad de las
respuestas y, por ende, la seguridad del sistema.

Diversas investigaciones han explorado variantes de PUFs basadas en
registros de desplazamiento y lógica reconfigurable sobre FPGAs.
Srilakshmi~et~al.~\cite{srilakshmi_2023_puf_lfsr} propusieron una PUF fuerte
construida con un LFSR de bajo consumo, orientada a dispositivos con
restricciones de energía. Hou~et~al.~\cite{hou_2019_lightweight} presentaron
una PUF LFSR ligera con seguridad mejorada implementada en FPGA, mientras que
Huang~et~al.~\cite{huang_2024_nlfsr} extendieron el concepto hacia registros
de desplazamiento no lineales (NLFSR PUF) para aumentar la resistencia frente
a ataques de modelado por aprendizaje automático.
Liu~et~al.~\cite{liu_2019_xor} exploraron PUFs reconfigurables basadas en
compuertas XOR como solución de bajo costo para seguridad en IoT. Las
revisiones exhaustivas de Lata y
Cenkeramaddi~\cite{lata_2023_fpgapuf} y de Anandakumar~et~al.
\cite{anandakumar_2021_fpgapufs} ofrecen un panorama amplio del estado del
arte en diseños de PUFs sobre FPGAs, abarcando desde arquitecturas basadas en
retardos hasta implementaciones basadas en memoria.

En el ámbito de la integración de PUFs dentro de sistemas embebidos y
protocolos de seguridad para IoT, Chen~et~al.~\cite{chen_2016_fpga}
implementaron generadores de números pseudoaleatorios criptográficamente
seguros basados en SRAM PUFs sobre FPGA.
Idriss~et~al.~\cite{idriss_2021_lightweight} propusieron un protocolo de
autenticación ligero basado en PUFs para dispositivos IoT con recursos
limitados, mientras que Hou~et~al.~\cite{hou_2023_modeling} desarrollaron un
protocolo de autenticación con doble PUF resistente a ataques de modelado
físico. Estos trabajos evidencian una tendencia creciente hacia la
incorporación de primitivas PUF como raíces de confianza en arquitecturas
embebidas.

No obstante, a pesar de los avances descritos, no se identificó en la
literatura un módulo PLPUF de código abierto, parametrizable y empaquetado
como núcleo IP compatible con la interfaz AXI4-Lite y el procesador
MicroBlaze~V. La mayoría de las implementaciones reportadas se limitan a
evaluaciones aisladas del núcleo PUF sin ofrecer integración dentro de un SoC
completo ni controladores de software reutilizables. El presente proyecto
aborda este vacío al proporcionar una solución integrada, desde el diseño HDL
hasta el software de evaluación, orientada a entornos académicos y de
investigación.

% ---------------------------------------------------------------------------
\section{Planteamiento del problema}

\subsection{Contexto del problema}

La implementación de sistemas seguros en plataformas reconfigurables, como las FPGAs,
se ha convertido en una necesidad crítica frente al incremento de amenazas relacionadas
con clonación de dispositivos, falsificación de hardware y accesos no autorizados.
Los mecanismos criptográficos convencionales dependen de claves almacenadas en memorias
no volátiles, tales como EEPROM, Flash o SRAM con respaldo de batería. Sin embargo,
estas tecnologías no garantizan protección frente a ataques físicos capaces de extraer
información sensible directamente del dispositivo \cite{hou_2019_lightweight}.

Esta situación resulta especialmente crítica en el contexto del IoT y sistemas embebidos
desplegados en entornos abiertos, donde los dispositivos operan con restricciones de
energía, área y capacidad computacional. La incorporación de módulos adicionales de
seguridad puede representar una sobrecarga significativa de recursos, dificultando la
adopción de soluciones robustas \cite{shamsoshoara_2020_puf, alhamarneh_2024_strengthening}.

Durante la fabricación de circuitos integrados surgen variaciones físicas inevitables,
tales como diferencias en el voltaje de umbral o el espesor del óxido de compuerta.
Aunque estas variaciones no afectan la funcionalidad lógica del circuito, constituyen
una fuente de entropía reproducible que puede explotarse para generar
identificadores únicos de hardware \cite{sudhanya_2016_silicon}.

\subsection{Justificación del problema}

Actualmente no se dispone de un módulo PUF de código abierto, programable y compatible
con el procesador suave MicroBlaze~V mediante la interfaz AXI. Muchas implementaciones
reportadas en la literatura presentan limitaciones de accesibilidad, flexibilidad o
integración dentro de arquitecturas modernas de sistemas embebidos, restringiendo su uso
en entornos académicos y de investigación \cite{lata_2023_fpgapuf, anandakumar_2021_fpgapufs}.

El desarrollo de un módulo PLPUF programable e integrable directamente en un SoC permite
eliminar la dependencia de almacenamiento persistente de secretos, mejorar la seguridad
a nivel físico y habilitar plataformas experimentales para la evaluación de técnicas
criptográficas emergentes \cite{delvaux_2015_helper, farha_2019_physical}. Asimismo,
facilita su adopción en contextos educativos y prototipos de investigación dentro del
ámbito nacional.

\subsection{Enunciado del problema}

Existe una vulnerabilidad crítica en sistemas embebidos basados en FPGA con procesadores
suaves debido a la dependencia de claves criptográficas almacenadas en memorias no
volátiles, las cuales pueden ser extraídas mediante ataques físicos. Esta vulnerabilidad
se ve agravada por la ausencia de módulos de seguridad de hardware abiertos,
programables y fácilmente integrables mediante estándares como AXI en arquitecturas
basadas en MicroBlaze~V dentro de FPGAs de la familia Artix-7.


% ---------------------------------------------------------------------------
\section{Objetivos del proyecto}

\subsection{Objetivo general}

Diseñar un módulo PLPUF programable compatible con el procesador MicroBlaze~V y la
interfaz AXI, mediante lenguajes de descripción de hardware y herramientas de síntesis
en FPGA, para su integración en sistemas embebidos orientados a la seguridad.

\subsection{Objetivos específicos}

\begin{enumerate}
\item Implementar el módulo PLPUF en HDL con ancho parametrizable que genere
identificadores únicos aprovechando variaciones físicas del hardware mediante
síntesis en FPGA con AMD Vivado.

\item Desarrollar una interfaz AXI-\textit{slave} que permita el control del módulo
desde el procesador MicroBlaze~V mediante registros de control y estado.

\item Integrar el módulo dentro de un sistema embebido completo con MicroBlaze~V,
memoria y UART, utilizando AMD Vitis para el desarrollo del software de control.

\item Evaluar el PLPUF mediante métricas de seguridad como uniformidad, unicidad y
confiabilidad utilizando pruebas automatizadas de desafío-respuesta y análisis
estadístico.
\end{enumerate}

% ---------------------------------------------------------------------------
\section{Alcances, entregables y limitaciones}

\subsection{Alcances}

El proyecto comprende el diseño, implementación e integración de un módulo PLPUF
programable en la placa Basys3 con FPGA Artix-7 (XC7A35T-1CPG236C). El módulo será
incorporado como periférico AXI-\textit{slave} dentro de un sistema SoC con procesador
MicroBlaze~V, memoria local y comunicación UART. Se desarrollará además el firmware de
control en lenguaje C utilizando AMD Vitis y se realizará la evaluación estadística del
módulo mediante métricas estándar de PUFs: uniformidad, unicidad y confiabilidad
\cite{maiti_2012_systematic, wilde_2020_confidence}.

\subsection{Entregables}

\begin{enumerate}
\item Código HDL del módulo PLPUF parametrizable con reportes de síntesis.
\item Módulo AXI-\textit{slave} funcional con documentación del mapa de registros.
\item Sistema SoC completo con diseño en Vivado, \textit{bitstream} y software de control.
\item \textit{Dataset} de pares desafío-respuesta e informe de evaluación estadística.
\end{enumerate}

\subsection{Limitaciones}

\begin{itemize}
\item El diseño está condicionado por los recursos disponibles en la FPGA Artix-7
asignada, lo que puede limitar parámetros como el ancho del PLPUF.

\item Las herramientas utilizadas serán exclusivamente AMD Vivado y AMD Vitis.

\item La evaluación estadística se realizará sobre un único dispositivo FPGA, por lo
que los resultados pueden no generalizarse a otros lotes de fabricación
\cite{wilde_2020_confidence}.
\end{itemize}

%%% Local Variables:
%%% mode: latex
%%% TeX-master: "main"
%%% End: